% Основной пресет настроек
\documentclass[a4paper,14pt,russian]{extreport}

% Основные пакеты
\usepackage{cmap}
\usepackage[T1, T2A]{fontenc}
\usepackage[utf8x]{inputenc}
\usepackage[russian]{babel}
\usepackage{listings}
\usepackage{graphicx}
\usepackage{setspace}
\usepackage{titlesec}

% Специальный пакет, чтобы изменить колонтикулы страниц
\usepackage{fancyhdr}
\pagestyle{fancy}
\fancyhf{}
\fancyhead[R]{\thepage}
\fancyheadoffset{0mm}
\fancyfootoffset{0mm}
\setlength{\headheight}{17pt}
\renewcommand{\headrulewidth}{0pt}
\renewcommand{\footrulewidth}{0pt}
\fancypagestyle{plain}{ 
    \fancyhf{}
    \rhead{\thepage}}
\setcounter{page}{1} % начать нумерацию страниц с №1

% Специальный пакет, чобы изменить стили
% всех заголовков в соответствии с ГОСТ 7.32-2001
% («Отчет о научно-исследовательской работе. Структура
% и правила оформления»
\usepackage{titlesec}
\titleformat{\chapter}[display]
    {\filcenter}
    {\MakeUppercase{\chaptertitlename} \thechapter}
    {8pt}
    {\bfseries}{}
 
\titleformat{\section}
    {\normalsize\bfseries}
    {\thesection}
    {1em}{}
 
\titleformat{\subsection}
    {\normalsize\bfseries}
    {\thesubsection}
    {1em}{}

% Полуторная высота строк
\renewcommand{\baselinestretch}{1.5}
 
% Настройки вертикальных и горизонтальных отступов
% для заголовков в соответствии с ГОСТ 7.32-2001
\titlespacing*{\chapter}{0pt}{-30pt}{8pt}
\titlespacing*{\section}{\parindent}{*4}{*4}
\titlespacing*{\subsection}{\parindent}{*4}{*4}

% Путь до директорий с картинками, которые
% включаются во все отчеты. Например логотип МГТУ.
\graphicspath{{/root/Learning/bmtsu-cpp/latex/images/}}

% Стили для листенгов
\lstset{
  frame=single, % Обрамлять все листинги рамкой
  breaklines=true % Переносить длинные строки
}

% Настройка паддингов страниц отчета в соответствии
% с ГОСТ 7.32-2001
\usepackage{geometry}
\geometry{left=3cm}
\geometry{right=1.5cm}
\geometry{top=2.4cm}
\geometry{bottom=2.4cm}

% Альтернативные настройки паддингов страниц
% \usepackage[left=30mm,right=20mm,top=20mm,bottom=20mm,headsep=0pt]{geometry}
\titleformat{\section}
  {\normalfont\fontsize{16}{16}\sffamily\bfseries}{\thesection}{1em}{}

% Устанавливаем Sans Serif
% \renewcommand{\familydefault}{\sfdefault}

%%%%%%%%%%%%%%%%%%%%%%%%%%%%%%%%%%%%%%%%%%%%%%%%%%%%%%%%
% Начало основной части документа
%%%%%%%%%%%%%%%%%%%%%%%%%%%%%%%%%%%%%%%%%%%%%%%%%%%%%%%%

\begin{document}

\centerline{\small{Федеральное государственное бюджетное образовательное учреждение}}
\centerline{\small{высшего профессионального образования}}

\vspace{1cm}

\begin{tabular}{ l l }
  \begin{tabular}[x]{@{}c@{}}\includegraphics[scale=0.8]{bmtsu}\end{tabular} & \begin{tabular}[x]{@{}c@{}}«Московский государственный технический\\университет имени Н.Э. Баумана»\\(МГТУ им. Н.Э. Баумана)\end{tabular} 
\end{tabular}

\vspace{1cm}

\begin{tabular}{ l l }
  ФАКУЛЬТЕТ & ИУ «Информатика и системы управления» \\
  КАФЕДРА & ИУ6 «Компьютерные системы и сети»
\end{tabular}

\vfill

\centerline{\huge{Отчет}}
\centerline{\large{по лабораторной работе №21}}

\vspace{0.5cm}
\centerline{\large{Дисциплина: ООП}}

\vspace{0.5cm}
\centerline{Название лабораторной работы: Тема 22. Использование указателей при работе со строками С++.}

% \vspace{3cm}
\vfill

\begin{tabular}{ l l }
  Выполнил студент группы ВИУ6-51 & Милкова В.Ю. \\
  Принял: & Ничушкина Т.Н.
\end{tabular}

\vfill

\centerline{Москва, \the\year}
\clearpage

\pagebreak



\section*{Задача}

Проверить правильность расположения в тексте круглых скобок (количество открывающихся и закрывающихся скобок должно совпадать). Пользуясь указателями определить адрес символов, расположенных между 2й открывающейся и ближайшей закрывающейся скобками.


\section*{Программный код}

\lstinputlisting[language=C++]{sources/main.cpp}


\section*{Результат работы}

\lstinputlisting{result/output.txt}


\section*{Блок схема}

\includegraphics[scale=0.5]{result/scheme.png}

\end{document}

